\section{Expresividad de la política y modelos generativos}

Un elemento clave del aprendizaje por imitación es cómo expresar la distribución de la política. Una regresión gaussiana única solo puede representar un comportamiento local y unimodal, y no puede igualar datos de expertos multimodales.

\subsection{Capacidad de representación de distribuciones}

Las acciones discretas pueden modelarse directamente con una distribución categorical, mientras que en el nivel de acciones continuas las formulaciones habituales también incluyen normalizing flow, Mixture Density Network y diffusion model.

\begin{importantbox}{El salto de expresividad que aportan los modelos generativos}
Al considerar la política como una red generativa $p_\theta(a \mid s)$, es posible retomar arquitecturas como diffusion, flow y VAE, de modo que la policy misma se convierta en un modelo probabilístico muestreable e inferible, lo que la adapta de manera natural a datos de demostración multimodales.
\end{importantbox}

\subsection{Políticas Diffusion, Mixture y Autoregressive}

Los medios de expresión más comunes:
\begin{enumerate}
    \item \textbf{Mixture of Gaussians}: produce varios $(\mu_i, \sigma_i, \alpha_i)$ a la vez y los combina con pesos softmax.
    \item \textbf{Autoregressive discretization}: divide el espacio de acciones en subdimensiones y modela de forma recursiva la distribución de cada dimensión.
    \item \textbf{Diffusion Policy}: primero muestrea ruido y después lo elimina de forma iterativa hasta generar la acción final.
\end{enumerate}

\begin{table}[H]
\centering
\begin{tabular}{p{0.22\textwidth} p{0.34\textwidth} p{0.34\textwidth}}
\toprule
Método & Ventaja & Escenario típico \\
\midrule
MoG & Ofrece un modo de mezcla limitado e interpretable & Control continuo de baja dimensión \\
Autoregressive & Descompone las dimensiones y captura correlation de grano fino & Secuencias de comportamiento de conducción autónoma \\
Diffusion & Permite implicit sampling & Manipulación robótica de alta dimensión \\
\bottomrule
\end{tabular}
\caption{Arquitecturas típicas de políticas expresivas}
\end{table}

\subsection{Muestreo y diagnóstico de cobertura}

Una vez entrenados los parámetros de la política, es necesario verificar su cobertura de la expert distribution. Las prácticas habituales incluyen:
\begin{itemize}
    \item KL divergence sweep: estima la KL con datos de expertos y datos muestreados.
    \item Reject sampling: si durante el rollout la probabilidad de la acción cae por debajo del threshold, se activa un human override.
    \item Latent space interpolation: interpola en las variables de control latent y observa si la acción generada cae dentro del manifold de los expertos.
\end{itemize}

\begin{knowledgebox}{Cambio de estrategia de muestreo}
Archit recuerda: \textit{«la estrategia de muestreo determina si podemos observar comportamientos extremos»}. Antes del despliegue, los ajustes de temperature y top-k/top-p de la policy deben escribirse en el runbook y registrarse en el observational log.
\end{knowledgebox}

Flow-based policy se usa para asignar el espacio de acciones a un latent space más regular; durante el entrenamiento se ajusta directamente la distribución objetivo en el latent, de modo que al muestrear basta con avanzar en la dirección latent y después recuperar la acción mediante un invertible map. Para acciones de alta dimensión (como el control humanoid), flow puede reducir considerablemente la dificultad del muestreo.

\begin{importantbox}{La intuición del control latente}
Flow policy equivale a hacer primero un muestreo gaussiano simple en el espacio latente y después convertirlo en una acción compleja mediante un invertible mapping. Esto permite trasladar el objetivo del aprendizaje por imitación al latent space, lo que simplifica la propagación del gradiente.
\end{importantbox}

\subsection{Resumen de la sección}

La expresividad de la política determina si el aprendizaje por imitación puede cubrir el expert space. Modelar la distribución multimodal con modelos generativos, junto con Flow-based latent control y rejection sampling, permite corregir de manera continua en el extremo del rollout, aproximarse a la expert distribution y ofrecer un punto de partida estable para el fine-tuning de downstream RL.

